\section{\S 2.2 分子结构与物质性质}\label{sect:分子结构与物质性质}

	\subsection{分子极性}\label{subsec:分子极性}

		\noindent\textbf{极性分子：}在电场中具有取向作用的分子、正电荷中心与负电荷中心\textcolor{red}{不重合}的分子。

		\noindent\textbf{非极性分子：}在电场中无取向作用的分子、正电荷中心与负电荷中心\textcolor{blue}{重合}的分子。

		\cemh{H2} 非极性键  \cemh{HCl} 极性键（\cemh{H^\delta+} 端与 \cemh{Cl^\delta-} 端）

		非极性分子 \cemh{CO2}（键的极性矢量和为 0） 极性分子 \cemh{H2O}（角形结构）

		\cemh{CCl4} 为正四面体形，键的极性矢量和为 0（非极性分子）

		影响分子极性的因素：
		\begin{itemize}
			\item 键的极性
			\item 分子空间构型
		\end{itemize}

		\begin{center}
		分子极性与键的极性、分子空间构型的关系
		\small
		\begin{tabular}{|c|c|c|c|}
			\hline
			\noindent\textbf{分子类型} & \textbf{空间构型} & \textbf{实例} & \textbf{分子极性}\\
			\hline
			单原子分子 & -- & 稀有气体 & 非极性 \\
			\hline
			\begin{tabular}[c]{@{}c@{}}双原子分子\\ \cemh{A2}\\ \cemh{AB}\end{tabular} & 直线形 & \begin{tabular}[c]{@{}c@{}}\cemh{H2}、\cemh{O2}、\cemh{N2}、\cemh{Cl2} 等\\ \cemh{HX}、\cemh{CO} 等\end{tabular} & \begin{tabular}[c]{@{}c@{}}非极性\\ 极性\end{tabular} \\
			\hline
			\begin{tabular}[c]{@{}c@{}}三原子分子\\ \cemh{ABA}\\ \cemh{ABA}\\ \cemh{ABC}\end{tabular} & \begin{tabular}[c]{@{}c@{}}直线形\\ 角形（V 形、折线形）\\ 角形或直线形\end{tabular} & \begin{tabular}[c]{@{}c@{}}\cemh{BeCl2}、\cemh{CO2}、\cemh{CS2} 等\\ \cemh{H2O}、\cemh{H2S}、\cemh{SO2} 等\\ \cemh{HCN}、\cemh{HClO}\end{tabular} & \begin{tabular}[c]{@{}c@{}}非极性\\ 极性\\ 极性\end{tabular} \\
			\hline
			\begin{tabular}[c]{@{}c@{}}四原子分子\\ \cemh{AB3}\\ \cemh{AB3}\end{tabular} & \begin{tabular}[c]{@{}c@{}}平面正三角形\\ 三角锥形\end{tabular} & \begin{tabular}[c]{@{}c@{}}\cemh{BF3}、\cemh{BCl3}、\cemh{SO3} 等\\ \cemh{NH3}、\cemh{PH3}、\cemh{PCl3}、\cemh{NF3} 等\end{tabular} & \begin{tabular}[c]{@{}c@{}}非极性\\ 极性\end{tabular} \\
			\hline
			\begin{tabular}[c]{@{}c@{}}五原子分子\\ \cemh{AB4}\\ \cemh{AB_xC_y}(x+y=4)\end{tabular} & \begin{tabular}[c]{@{}c@{}}正四面体形\\ 四面体形\end{tabular} & \begin{tabular}[c]{@{}c@{}}\cemh{CH4}、\cemh{CCl4} 等\\ \cemh{CH3Cl}、\cemh{CH2Cl2}、\cemh{CHCl3}\end{tabular} & \begin{tabular}[c]{@{}c@{}}非极性\\ 极性\end{tabular} \\
			\hline
		\end{tabular}
		\end{center}

		\begin{center}
		\begin{minipage}[t]{0.56\textwidth}
		\raggedright
		\noindent\textbf{基本有机分子：}
		\begin{tabular}{@{}l l@{}}
			乙炔： & \makebox[4em]{\cemh{HC#CH}}\quad 直线形\quad 非极性 \\
			乙烯： & \makebox[4em]{\cemh{H2C=CH2}}\quad 平面形\quad 非极性 \\
			苯： & \makebox[4em]{\chemfig[atom sep=1em]{*6(-=-=-=)}}\quad 平面形\quad 非极性
		\end{tabular}
		\end{minipage}
		\hfill
		\begin{minipage}[t]{0.36\textwidth}
		\raggedright
		\noindent\textbf{特殊分子：}

		\begin{tabular}{@{}l l@{}}
			\cemh{H2O2} & 极性 \\
			\cemh{N2H4} & 极性
		\end{tabular}
		\end{minipage}
		\end{center}

		\noindent\textcolor{blue}{【归纳】}
		\begin{enumerate}
			\item 仅由非极性键形成的分子一定为非极性分子（键无极性，分子无极性）。
			\item 由极性键形成的双原子分子一定为极性分子。
			\item 由极性键形成的多原子分子，若空间结构有对称中心，则为非极性分子；若无对称中心，则为极性分子。
			\item 对于 \cemh{AB_n} 分子：若 A 原子价层电子都已成键（都是$n+0$型），则 \cemh{AB_n} 为非极性分子\textcolor{red}{（可由 VSEPR 证明）}；若 A 原子价层电子未都成键，则 \cemh{AB_n} 为极性分子。
		\end{enumerate}

	\subsection{相似相溶原理}\label{subsec:相似相溶原理}

		\cemh{HCl}、\cemh{HBr}、\cemh{NH3} 极易溶于水；\cemh{Br2}、\cemh{I2} 易溶于 \cemh{CCl4}、苯等有机溶剂。

		\textcolor{blue}{极性分子易溶于极性溶剂，非极性分子易溶于非极性溶剂。}

		\textcolor{red}{相似相溶原理——经验规律}

		\noindent\textcolor{red}{【问】物质普遍存在三态变化，由分子构成的物质三态变化时也伴随能量变化，这说明什么？}

		存在分子间作用力。

	\subsection{范德华力}\label{subsec:范德华力}

		\noindent\textbf{范德华力：}分子之间普遍存在的相互作用力，是最常见的分子间作用力。

		比化学键弱得多：几\si{\kilo\joule\per\mole}）

		\textcolor{red}{\(\therefore\)分子构成的物质通常熔沸点较低、硬度较小}

		本质上也是电性相互作用，无方向性和饱和性，随分子间距离增大而迅速减弱

		影响范德华力大小的因素：
		\begin{itemize}
			\item 分子组成和结构相似时：相对分子质量越大，范德华力越强。
			\item 结构相似、相对分子质量接近时：分子极性越强，范德华力通常越强。
		\end{itemize}

		\begin{xmp}{沸点异常现象}
			比较 \cemh{HF}、\cemh{HCl}、\cemh{HBr}、\cemh{HI} 的沸点高低；比较 \cemh{H2O}、\cemh{H2S}、\cemh{H2Se}、\cemh{H2Te} 的沸点高低

			\textcolor{blue}{\cemh{NH3}、\cemh{H2O}、\cemh{HF} 存在沸点反常升高现象}
			
			\textcolor{red}{说明除了范德华力外还存在更强的分子间作用力}
		\end{xmp}

	\subsection{氢键}\label{subsec:氢键}

		\noindent\textbf{定义：}一个分子中电负性很强的原子与另一个分子中的氢原子之间形成的较强相互作用。

		\textcolor{red}{氢键的本质：（分子间的）静电作用}

		\noindent\textbf{表示方式：}\cemh{X-H\bond{...}Y}（\cemh{X}、\cemh{Y} 可相同，也可不同）。

		\noindent\textbf{形成条件：}\cemh{H} 原子与 \cemh{N}、\cemh{O}、\cemh{F} 相连。

		\textcolor{red}{氢键作用能：几十\si{\kilo\joule\per\mole}}\textcolor{blue}{（比范德华力强、比化学键弱）}

		\subsubsection{氢键对物质性质的影响}\label{subsubsec:氢键对物质性质的影响}

			\begin{enumerate}
				\item \cemh{H2O}、\cemh{HF}、\cemh{NH3} 的熔沸点高于同族其他氢化物
				\item 水的比热容较大\textcolor{blue}{（缔合水分子：\cemh{(H2O)_n}）}
				\item 冰的密度小于液态水\textcolor{red}{（氢键有方向性和饱和性）}
				
				\textcolor{red}{由于氢键作用，冰中的水分子整齐规则排列，分子间隙大于液态水，所以密度减小}

				\item 氨易液化且极易溶于水（\cemh{NH3} 分子间、\cemh{NH3} 与 \cemh{H2O} 间形成氢键）
				\item 乙醇与水互溶
				
				\chemfig[atom sep=2.2em]{
					C_2H_5-O(-[:30,,,,dotted]H-[:30]O-[:330]H)-[:330]H-[:330,,,,dotted]O(-[:210]H)-[:330]H
				}

				\item 羧酸沸点较高
				
				\textcolor{blue}{双分子缔合}：
				\chemfig[atom sep=2.2em]{CH_3-C=[:60]O-[:0,,,,dotted]H-O-[:300]C(-CH_3)=[:240]O-[:180,,,,dotted]H-[:180]O-[:120]\phantom{C}}
				
				\item 分子内氢键与分子间氢键	
					\begin{center}
					\small
					\begin{tabular}{|c|c|c|}
						\hline
						\noindent\textbf{物质} & \textbf{熔点/\si{\celsius}} & \textbf{沸点/\si{\celsius}} \\
						\hline
						邻羟基苯甲醛（分子内氢键） & 2 & 195.5 \\
						\hline
						对羟基苯甲醛（分子间氢键） & 115 & 246.6 \\
						\hline
						邻苯二酚 & 103 & 245 \\
						\hline
						对苯二酚 & 173 & 286 \\
						\hline
					\end{tabular}
					\end{center}

					\textcolor{blue}{组成结构相似时，含分子内氢键的物质熔沸点通常低于含分子间氢键的物质。}
				\item 蛋白质二级结构：\chemfig[atom sep=2.2em]{-C(=[2]O)-}与\chemfig[atom sep=2.2em]{-N(-[2]H)-}之间形成氢键
				\item DNA 双螺旋中碱基相互配对
			\end{enumerate}
